\section{The strict Selmer comparison and its numerical bound}
\label{sc-section}

The arithmetic complex in Section~\ref{ac-section} proves the
homology torsion estimate used by the submitted class-counting
argument.  It does not by itself bound a strict Selmer group.
We now state the comparison assertions precisely,
prove one elementary improvement in their numerical factor,
and give the resulting conditional exponent.
This section isolates the comparisons before proving the numerical
consequence. Sections~\ref{sl-section}, \ref{su-section} and \ref{sp-section} then prove
the class-number estimate and assertions~\ref{sc-count}--\ref{sc-layer}.
Theorem~\ref{sw-few} later proves the finite-field homology
comparison in assertion~\ref{sc-few}, with an improved scalar exponent.

\subsection{The precise arithmetic objects}

Let $\ell$ be prime and let $E\ge1$.  For a number field $L$
put $n_L=[L:\mathbb Q]$, $D_L=|\operatorname{disc}L|$,
$r_2(L)=\#\{\text{complex places of }L\}$, and
\[
 t_L(N,m)=\#H_m(\operatorname{GL}_N(\mathcal O_L),
                       \mathbb Z)_{\rm tors},
 \qquad
 g_L=2n_L+\lfloor\log_2 h_L\rfloor+2,
 \tag{SC1}
\]
where $h_L$ is the ideal class number.
For the totally imaginary fields occurring below, the relevant
homology torsion groups are finite by Section~\ref{ac-section}.

Let $\mathcal S_\ell(L,E)$ denote the strict Selmer group with
coefficients $\mathbb Q_\ell/\mathbb Z_\ell(-E)$: it consists
of the classes in $H^1(G_L,\mathbb Q_\ell/\mathbb Z_\ell(-E))$
which are unramified outside $\ell$ and have zero restriction
to every decomposition group above $\ell$.
Equivalently it is the kernel of localization above $\ell$
in $H^1(G_{L,S},\mathbb Q_\ell/\mathbb Z_\ell(-E))$,
where $S$ is the set of primes above $\ell$.
Write
\[
       s_\ell(L,E)=
       \length_{\mathbb Z_\ell}\mathcal S_\ell(L,E)
             \in\mathbb N\cup\{\infty\}.
       \tag{SC2}
\]
When this length is finite, the group has order
$\ell^{s_\ell(L,E)}$.
These are the negative twists used by the source.

Consider a finite cyclic Galois extension $K/L$ with generator
$\gamma$, and integers $\nu\ge c\ge1$, with $c\ge2$ if $\ell=2$.
Suppose that $K$ contains a primitive $\ell^\nu$-th root
$\zeta$ satisfying
\[
          \gamma(\zeta)=\zeta^{\,1+\ell^c},
          \qquad [K:L]=\ell^{\nu-c}.
          \tag{SC3}
\]
Put $q=\ell^\nu$, $\kappa=1+\ell^c$, and $a=\kappa^E$.
For a fractional divisor $\mathfrak d=\sum_P d_P P$
on $K$, call $\alpha\in K^\times$ admissible for
$\mathfrak d$ if
\[
                v_P(\alpha)=d_{\gamma P}-a d_P
                       \quad\text{for every finite prime }P.
                \tag{SC4}
\]
Define the explicitly weighted product
\[
       N_{\gamma,a}(\alpha)
          =\prod_{j=0}^{[K:L]-1}\gamma^j(\alpha)^{a^j}.
          \tag{SC5}
\]
This is not the ordinary field norm.

Let $T_{\gamma,a,q}$ be the set of functions
$\varphi:\{P\text{ finite in }K\}\to\mathbb Z/q\mathbb Z$
such that
\[
 \sum_P v_P(\alpha)\varphi(P)=0\quad(\alpha\in K^\times),
 \qquad
 \varphi(\gamma^{-1}P)=a\varphi(P)\quad(P).
 \tag{SC6}
\]
The sum is finite because a principal divisor has finite
support.  The first condition identifies these functions
with homomorphisms out of the ideal class group; the second
is the displayed equivariance condition.  In particular
the inverse on $\gamma$ is part of the statement.

For $b,B\in\mathbb N$, write $\mathcal A(b,B)$ for the
following assertion about these fixed data.  Every finite
set $S\subset K^\times$ has cardinality at most $B$
provided that each $\alpha\in S$ is admissible for some
fractional divisor and that, for all $\alpha,\alpha'\in S$,
\[
 \begin{split}
 N_{\gamma,a}(\alpha)^{\ell^b}
     &=N_{\gamma,a}(\alpha')^{\ell^b}\delta^q
       \text{ for some }\delta\in K^\times\\
 &\hspace{35mm}\Longrightarrow\quad\alpha=\alpha'.
 \end{split}
 \tag{SC7}
\]
Thus this assertion bounds a set of distinct classes
represented by the weighted products, after the specified
power is taken.  It does not assert a bound on the infinite
set of all admissible elements.

\subsection{The four comparison assertions}

The following are the exact mathematical assertions used in this
route of the submitted proof\cite{alpo-uniform-source}.

\begin{enumerate}
\item\label{sc-few}
For totally imaginary $L$ and data (SC3), assume
\[
                 2E+1<\ell^c,\qquad 2E+1<r_2(L).
                 \tag{SC8}
\]
Then $\mathcal A(b,B)$ holds with
\[
 \begin{split}
 b&=2E+(2E+1)(c+2E+1),\\
 B&=q^{r_2(L)}t_L(4E+6,2E+1)t_L(4E+6,2E).
 \end{split}
 \tag{SC9}
\]
In the source this assertion is obtained from constructed
finite-field homology functionals, integral homology
classes with prescribed reductions, an odd-degree rational
rank estimate, and detection of global powers by residue
fields.  Those constructions and comparisons are not
replaced here by the homology torsion bound.

\item\label{sc-count}
For totally imaginary $L$, any data (SC3), and any $b,B$
for which $\mathcal A(b,B)$ holds, one has
\[
 |T_{\gamma,a,q}|\,q^{r_2(L)}
       \le B\,\ell^{(b+v_\ell(E))g_L+c+v_\ell(E)}.
 \tag{SC10}
\]
This is the separate comparison with class groups and
units.  It does not require the two inequalities (SC8).

\item\label{sc-sha}
Suppose $L$ is totally imaginary and $c\ge1$, with $c\ge2$
at $\ell=2$, is the largest integer such that
$\mu_{\ell^c}\subset L$.  If there is a single integer $B$
such that
\[
                         |T_{\gamma,\kappa^E,\ell^\nu}|
                              \le B
                              \tag{SC11}
\]
for every $\nu$ and every datum (SC3) over $L$, then
$s_\ell(L,E)$ is finite and
\[
                              \ell^{s_\ell(L,E)}\le B.
                              \tag{SC12}
\]
Uniformity in the entire cyclotomic tower is essential.

\item\label{sc-layer}
Suppose $L$ contains $\mu_\ell$ if $\ell$ is odd and
$\mu_4$ if $\ell=2$.  There is a totally imaginary
cyclotomic extension $L'/L$ of degree $d$, and a largest
integer $c$ with $\mu_{\ell^c}\subset L'$, such that
\[
 \begin{gathered}
 c\ge1,\quad \ell=2\Longrightarrow c\ge2,\qquad
 2E+1<\ell^c,\quad 2E+1<r_2(L'),\\
 n_{L'}=dn_L,\quad
 d\le(n_L+1)(4E+4),\quad
 \ell^c\le2n_{L'},\\
 D_{L'}\le D_L^d(\ell^c)^{n_{L'}},\qquad
 s_\ell(L,E)\le s_\ell(L',E).
 \end{gathered}
 \tag{SC13}
\]
The numerical statement alone does not prove the last
inequality; that is the Selmer comparison in this assertion.
\end{enumerate}

For the numerical conclusion below we also isolate the
classical class-number bound used by the source:
\[
 \log h_L\le\tfrac12\log D_L+
             n_L\log(1+\tfrac12\log D_L).
 \tag{SC14}
\]
Its source proof invokes the existence of a small ideal
in every ideal class and a count of ideals of bounded norm.
Those inputs are not proved in this section.
We will only need the weaker consequence
\[
                        g_L\le2n_L+4n_L\log(2D_L)+2.
                        \tag{SC15}
\]
Indeed $\log(1+x)\le x$ gives
$\log h_L\le n_L\log D_L$, since $n_L\ge1$.
The inequality $\log2>1/2$ then gives
$\lfloor\log_2h_L\rfloor\le2n_L\log D_L
\le4n_L\log(2D_L)$.

\subsection{An exact valuation in the class construction}

The large value of $b$ in (SC9) comes partly from bounding
a product of differences of powers.  That product admits
a simple exact evaluation of its valuation.

\begin{lemma}\label{sc-vandermonde}
Let $\ell$ be prime, $c\ge1$, and $c\ge2$ if $\ell=2$.
Put $\kappa=1+\ell^c$.  For $E\ge1$ define
\[
             P_E=\prod_{\substack{0\le j\le2E\\j\ne E}}
                       (\kappa^E-\kappa^j)\in\mathbb Z.
\]
Then $P_E\ne0$ and
\[
                   v_\ell(P_E)=2Ec+2v_\ell(E!).
                   \tag{SC16}
\]
Consequently, in any $\mathbb Z/\ell^\nu\mathbb Z$-module,
\[
           \ell^{2E}P_E z=0
             \quad\Longrightarrow\quad
           \ell^{\,2E(c+1)+2v_\ell(E!)}z=0.
           \tag{SC17}
\]
\end{lemma}
\begin{proof}
We first establish the valuation identity
\[
                   v_\ell(\kappa^m-1)=c+v_\ell(m)
                              \quad(m\ge1).
                   \tag{SC18}
\]
If $t$ has valuation $h\ge1$ (and $h\ge2$ for $\ell=2$),
and $u$ is a positive integer prime to $\ell$, the binomial
expansion of $(1+t)^u-1$ has first term $ut$ of valuation
$h$ and all later terms of valuation at least $2h>h$.
If instead the exponent is $\ell$, its first term
$\ell t$ has valuation $h+1$.  All terms with
$2\le j<\ell$ have valuation at least $1+2h>h+1$,
because $\ell\mid\binom{\ell}{j}$.
The last term $t^\ell$ also has strictly greater valuation:
for odd $\ell$, $\ell h\ge h+2$; for $\ell=2$ this uses
$h\ge2$.  Thus raising $1+t$ to the $\ell$-th power
increases the valuation of its difference from one by
exactly one.  Write $m=\ell^r u$ with $\ell\nmid u$,
first apply the prime-to-$\ell$ calculation, and then
apply the second calculation $r$ times.  This proves
(SC18).

The powers of $\kappa>1$ are distinct, so $P_E\ne0$.
For $j\ne E$ the unit factor
$\kappa^{\min(j,E)}$ and (SC18) give
\[
             v_\ell(\kappa^E-\kappa^j)
                          =c+v_\ell(|E-j|).
\]
The distances $|E-j|$ run through $1,\ldots,E$ twice.
Adding the valuations proves (SC16).
Write $P_E=\ell^\delta u$ with
$\delta=2Ec+2v_\ell(E!)$ and $\ell\nmid u$.
Multiplication by $u$ is invertible modulo $\ell^\nu$.
Cancelling it proves (SC17), including the case
$\delta\ge\nu$.
\end{proof}

The source uses precisely $P_E$ at this stage, with the
index in its definition shifted by $e=E-1$.
Its much larger valuation allowance is
$(2E+1)(c+2E+1)$.  Formula (SC17) supplies the smaller
power
\[
                         b_{\rm new}=2E(c+1)+2v_\ell(E!)
                         \tag{SC19}
\]
in the same finite-field detection calculation.
This proves the required cancellation of the scalar.
It does not prove the other constructions entering
assertion~\ref{sc-few}, nor does assertion~\ref{sc-few}
with its original exponent alone imply its stronger
version with $b_{\rm new}$.  Theorem~\ref{sw-few} carries
out those constructions with this scalar estimate and proves
the strengthened version directly.
For example, at $\ell=2$, $c=2$, $E=1$, the original
exponent is $17$, whereas (SC19) is $6$.

The restriction at $2$ cannot be dropped.
For $\ell=2$, $c=1$, and $m=2$ one has
$v_2(3^2-1)=3$, while $c+v_2(m)=2$.
The source includes $c\ge2$ in this case, so this is
a counterexample to the weakened statement only.

\subsection{The conditional exponent}

\begin{proposition}\label{sc-exponent}
Assume assertions~\ref{sc-few}--\ref{sc-layer} and the
bound (SC15).  Suppose that, for some real $C$ and every
totally imaginary number field $F$,
\[
 \log t_F(4E+6,m)
      \le\exp\bigl((2E n_F)^C\log(2D_F)\bigr)
       \qquad(0\le m\le2E+1).
 \tag{SC20}
\]
If $L$ contains the roots of unity required in
assertion~\ref{sc-layer}, then $s_\ell(L,E)$ is finite and
\[
 s_\ell(L,E)\log\ell
   \le\exp\bigl((2E n_L)^{\,5(\max(C,9)+2)+11}
                          \log(2D_L)\bigr).
 \tag{SC21}
\]
In particular the proved homology exponent $C=13$ gives
the conditional strict Selmer exponent $86$.
\end{proposition}
\begin{proof}
Choose $L',c,d$ from (SC13), and write
$n'=n_{L'}$, $D'=D_{L'}$, $r'=r_2(L')$, and
\[
 T_1=t_{L'}(4E+6,2E+1),\quad
 T_2=t_{L'}(4E+6,2E),\quad
 b=2E+(2E+1)(c+2E+1),\quad
 v=v_\ell(E).
\]
Apply assertions~\ref{sc-few} and \ref{sc-count} to
each datum (SC3) over $L'$.  The positive factor $q^{r'}$
cancels, leaving the bound, uniform in $\nu$,
\[
                |T_{\gamma,\kappa^E,\ell^\nu}|
                         \le T_1T_2\ell^X,
              \qquad X=(b+v)g_{L'}+c+v.
 \tag{SC22}
\]
Assertion~\ref{sc-sha} gives a finite length $s'=s_\ell(L',E)$
and $\ell^{s'}\le T_1T_2\ell^X$.  Hence
\[
           s'\log\ell\le\log T_1+\log T_2+X\log\ell.
           \tag{SC23}
\]

Put $u=2En'\ge2$ and $Y=\log(2D')>1/2$.
The inequalities in (SC13) give
\[
 \log\ell\le u,\quad c\le u,\quad v\le u,\quad
 b\le u+2u^2,\quad
 g_{L'}\le u+2uY+2.
\]
For the bound on $b$, use $2E\le u$,
$2E+1\le n'$, $c\le2n'$, and $2n'\le u$.
It follows that
$b+v\le3u^2$, $g_{L'}\le6uY$, and $c+v\le4uY$.
Consequently
\[
                  X\log\ell\le19u^4Y\le u^9Y,
                  \tag{SC24}
\]
since $u^5\ge32>19$.
By (SC20), (SC23) is at most
$2\exp(u^CY)+u^9Y$.

Let $A=u^{\max(C,9)}Y\ge1$.  Then
\[
 2\exp(u^CY)+u^9Y
       \le2e^A+A\le3e^A\le e^{4A}
       \le\exp\bigl(u^{\max(C,9)+2}Y\bigr).
 \tag{SC25}
\]
Here $A\le e^A$, $3\le e^{3A}$, and $u^2\ge4$
justify the three last comparisons.

It remains to express the bound in the original field.
Put $x=2E n_L\ge2$, $Y_0=\log(2D_L)>1/2$,
$w=\ell^c$, and $C_1=\max(C,9)+2\ge0$.
The extension bounds imply
\[
 d\le16E n_L=8x,\quad
 n'\le16E n_L^2\le4x^2,\quad
 u=2En'\le8x^2\le x^5,\quad w\le8x^2.
\]
Using the discriminant inequality in (SC13) yields
\[
 \begin{split}
 Y&\le\log2+d\log D_L+n'\log w\\
  &\le(1+8x)Y_0+32x^4
   \le73x^4Y_0
   \le x^{11}Y_0,
 \end{split}
\]
where $\log w\le w$, $Y_0>1/2$, and $x^7\ge128>73$
were used.  Thus $u^{C_1}Y\le x^{5C_1+11}Y_0$.
Combining this with (SC25) and
$s_\ell(L,E)\le s'$ proves (SC21).
The value for $C=13$ is $5(15)+11=86$.
\end{proof}

The original source inserts $C=24$ and obtains $141$
by this same exponent propagation.  The improvement to
$86$ follows once the comparison assertions are supplied.
The following sections prove the class-number input and the
other comparisons, and Theorem~\ref{sw-few} proves
assertion~\ref{sc-few}; Corollary~\ref{sw-coarse-consequence}
records the resulting bound.
It is not a proof of their conclusion from the geometric
complex alone.



